\subsection{Importance of Re-testing Updated Versions}\label{sec:needtesting}

When software developers release an updated version of a program, re-testing updated program is important because the program bugs may still remain after the update. This need for re-testing has been emphasized in prior researches~\cite{wang2024efficient,xiang2024critical,yi2022feedback,guo2016conc}. To quantify the impact of program bugs in version updates, we collected 1,022 release notes from GNU C programs commonly used in symbolic execution, including Coreutils~\cite{gnu_coreutils}, and counted how many release notes describe each factor as the reason for the update. 

Figure~\ref{fig:motivation} shows that bugs discovered in previous versions are the largest factor for program updates. Among the 1,022 release notes, 80.1\% of them mention that bugs discovered in the previous version were fixed. The result indicates that new program bugs are continuously discovered in updated program versions even after bugs in previous versions have been fixed. This is due to two main reasons: (1) vulnerabilities in newly added code and (2) conflicts between bug-fixing changes and other parts of the code. Thus, it is important to re-test the modified code to detect additional program bugs before releasing a new version. 